\documentclass[11pt]{article}
\usepackage[margin=1in]{geometry}
\usepackage{xcolor}
\usepackage{titlesec}
\usepackage{booktabs}
\usepackage{tabularx}
\usepackage{enumitem}
\usepackage[hidelinks]{hyperref}

\definecolor{accent}{HTML}{0B6E4F}
\titleformat{\section}{\large\bfseries\color{accent}}{}{0em}{}
\titlespacing*{\section}{0pt}{1.1em}{0.4em}
\setlength{\parindent}{0pt}
\setlength{\parskip}{0.55em}
\newcolumntype{R}{>{\raggedleft\arraybackslash}X}

\begin{document}

\begin{center}
  {\LARGE\bfseries\color{accent} Foundation for Open Science}\\[0.2em]
  {\large Research Grant Proposal}
\end{center}

\vspace{0.6em}
\begin{tabularx}{\textwidth}{@{}lX@{}}
\toprule
\textbf{Project title} & Mapping Hippocampal Replay to Consolidated Memory in Aging \\
\textbf{Principal investigator} & Dr.\ Priya Nadeau, Center for Neural Dynamics \\
\textbf{Co-investigators} & Dr.\ Samuel Okoro, Dr.\ Lena Fischer \\
\textbf{Host institution} & Aldermere University \\
\textbf{Requested amount} & \$742,000 over 3 years \\
\textbf{Program} & Cognitive Neuroscience of Aging (Cycle 12) \\
\bottomrule
\end{tabularx}

\section{Abstract}
Sharp-wave ripple events in the hippocampus replay recent experience and are
thought to drive the transfer of memories to neocortex during rest. Whether
this replay machinery degrades with normal aging, and whether that decline
predicts memory loss, remains unresolved. We combine high-density
electrophysiology in behaving rodents with a computational model of replay
scheduling to test the hypothesis that reduced replay fidelity, not reduced
replay rate, underlies age-related consolidation deficits.

\section{Specific Aims}
\begin{itemize}[leftmargin=1.4em, itemsep=0.15em]
  \item \textbf{Aim 1.} Quantify replay fidelity across the lifespan using
        Bayesian decoding of place-cell sequences during rest.
  \item \textbf{Aim 2.} Test whether closed-loop disruption of low-fidelity
        ripples worsens consolidation while sparing encoding.
  \item \textbf{Aim 3.} Fit a generative model that links single-ripple content
        to next-day recall performance.
\end{itemize}

\section{Budget Summary}
\begin{tabularx}{\textwidth}{@{}lRRRR@{}}
\toprule
\textbf{Category} & \textbf{Year 1} & \textbf{Year 2} & \textbf{Year 3} & \textbf{Total} \\
\midrule
Personnel        & 118{,}000 & 122{,}000 & 126{,}000 & 366{,}000 \\
Equipment        &  84{,}000 &  12{,}000 &  10{,}000 & 106{,}000 \\
Animals \& supplies &  46{,}000 &  48{,}000 &  49{,}000 & 143{,}000 \\
Travel \& dissemination &  9{,}000 &  9{,}000 &  11{,}000 &  29{,}000 \\
Indirect costs   &  32{,}000 &  33{,}000 &  33{,}000 &  98{,}000 \\
\midrule
\textbf{Total}   & \textbf{289{,}000} & \textbf{224{,}000} & \textbf{229{,}000} & \textbf{742{,}000} \\
\bottomrule
\end{tabularx}

\section{Broader Impacts}
The project trains two doctoral students in closed-loop neurophysiology, releases
all decoding software under an open license, and partners with a community
science program to bring memory research into secondary school classrooms.

\end{document}
